% ============================================================
% SECCIÓN 2
% ============================================================
\section{Masas y Pesos Sísmicos}
\label{sec:masas}

El peso sísmico del edificio se determina a partir de las cargas actuantes en el 
modelo, conforme a lo establecido en NCh433~\cite{NCh433:1996}. Para ello, en el 
apartado \textit{Mass Source} de ETABS se asignan los siguientes factores de 
combinación:

\begin{itemize}
    \item Peso propio (PP): factor $1{,}0$
    \item Sobrecargas (SC): factor $0{,}25$
\end{itemize}

Sumando la masa de todos los pisos y multiplicandolos por la aceleración de gravedad, se obtiene el peso sísmico de la estructura:
\begin{equation}
    P = 7{.}660{,}05\ \text{tonf}
\end{equation}

Comparado con el valor de prediseño ($P_{\text{pred}} = 8{.}766{,}65\ \text{tonf}$), 
la diferencia es de un \textbf{12.6\%}, valor que se encuentra dentro de los márgenes 
esperados para esta etapa del análisis.